%-------------------------
% JPL Scientist II cover letter for Nhat-Minh Nguyen
% Adapted from an entry-level LaTeX cover-letter template.
%------------------------

\documentclass[11pt,letterpaper]{article}

\usepackage{latexsym}
\usepackage[empty]{fullpage}
\usepackage{titlesec}
\usepackage{marvosym}
\usepackage[usenames,dvipsnames]{color}
\usepackage{verbatim}
\usepackage[hidelinks]{hyperref}
\usepackage{fancyhdr}
\usepackage{multicol}
\usepackage{csquotes}
\usepackage{tabularx}
\hypersetup{colorlinks=true,urlcolor=black}
\usepackage[11pt]{moresize}
\usepackage{setspace}
\usepackage{fontspec}
\usepackage[inline]{enumitem}
\usepackage{ragged2e}
\usepackage{anyfontsize}

\usepackage[letterpaper,margin=0.68in]{geometry}
\setlength{\footskip}{8pt}

\setmainfont[
BoldFont=SourceSansPro-Semibold.otf,
ItalicFont=SourceSansPro-RegularIt.otf
]{SourceSansPro-Regular.otf}

\pagestyle{fancy}
\fancyhf{}
\fancyfoot{}
\renewcommand{\headrulewidth}{0pt}
\renewcommand{\footrulewidth}{0pt}

\urlstyle{same}
\raggedbottom
\raggedright
\setlength{\tabcolsep}{0in}

\definecolor{UI_blue}{RGB}{32, 64, 151}
\definecolor{HF_color}{RGB}{179, 179, 179}

\titleformat{\section}{
\color{UI_blue} \scshape \raggedright \large
}{}{0em}{}[\vspace{-0.7cm} \hrulefill \vspace{-0.2cm}]

\begin{document}

\begin{center}
    {\Huge Nhat-Minh Nguyen} \\
    \vspace{0.06cm}
    {\color{UI_blue} \Large Scientist II, Roman and Euclid Cosmology} \\
    \vspace{0.08cm}
    \large
    \href{mailto:nhat.minh.nguyen@ipmu.jp}{nhat.minh.nguyen@ipmu.jp}
    \quad | \quad
    \href{https://linkedin.com/in/minhmpa}{linkedin.com/in/minhmpa}
    \quad | \quad
    \href{https://github.com/MinhMPA}{github.com/MinhMPA} \quad | \quad \href{https://orcid.org/0000-0002-2542-7233}{ORCID}

\vspace{0.05cm}
{\color{UI_blue} \hrulefill}
\end{center}

\justify
\setlength{\parindent}{0pt}
\setlength{\parskip}{8.5pt}
\vspace{-0.05cm}

\today

Search Committee \\
Jet Propulsion Laboratory

Dear Search Committee,

I am applying for the Scientist II position in Roman and Euclid cosmology because it is the kind of role my research has been preparing me for: helping turn precision maps of the late Universe into reliable tests of the dark sector. I work on large-scale-structure modeling, field-level Bayesian inference, and galaxy surveys. My goal is simple. Use as much information as we can from cosmic structure without letting modeling choices outrun the physics.

Roman and Euclid will make that balance difficult in the right way. Euclid brings wide-area coverage. Roman brings depth and image quality. Their overlap can test whether weak lensing and galaxy clustering are telling the same story, especially once nonlinear structure, galaxy bias, intrinsic alignments, photometric redshifts, selection effects, and cross-probe covariance all matter at the same time. These are not just technical details. They decide whether small error bars become reliable statements about dark energy, dark matter, and gravity.

Two parts of my record are especially relevant. First, my first-author work reported evidence for suppressed late-time structure growth, using information from the cosmic microwave background, weak lensing, galaxy clustering, and cosmic velocities to test the consistency of the standard cosmological model. Second, I developed field-level Bayesian inference methods for galaxy surveys, work recognized by the 2024 Buchalter Cosmology Prize. The projects are different in detail, but they come from the same question: what can we learn from the cosmic web if we model it carefully enough?

As a Scientist II, I would focus on careful technical work inside Roman and Euclid analysis teams. I would be especially interested in joint weak-lensing and clustering modeling, simulation-based likelihood checks, Roman-Euclid overlap tests, and dark-sector interpretation through growth, gravity, and dark energy. I do not think every analysis should immediately become field-level. The value of field-level inference is that it gives us a controlled way to ask what information is lost, where summary statistics are safe, and which assumptions start to fail as the uncertainties shrink.

JPL is the place for this work because Roman and Euclid require nonlinear structure modeling, statistical inference, and mission-scale collaboration to meet in the same analysis. I have worked with PFS and DESI, developed clustering and multi-survey calibration projects, contributed to community validation efforts, and mentored early-career researchers on inference and reconstruction problems. I would approach this role with focus and care. I would welcome the opportunity to contribute to JPL Roman and Euclid cosmology program.

\vspace{0.18cm}
Sincerely, \\
Nhat-Minh Nguyen

\end{document}
